\begin{frame}{2차원 투영: 동물 vs 과일 두 군집}
  \begin{columns}[T]
    \begin{column}{0.46\textwidth}
      \centering
      \includegraphics[width=\linewidth]{ch14_2_word2vec_2d.png}
      \vskip0.3em
      {\footnotesize 10개 단어 임베딩을 8차원 $\to$ 2차원(SVD 상위 2
      좌표)으로 투영.}
    \end{column}
    \begin{column}{0.52\textwidth}
      \begin{itemize}
        \item 10개 단어가 ``동물''/``과일'' 두 군집으로 \textbf{거의}
          나뉨 --- 왼쪽: 고양이는·강아지는·사자는·동물이다·귀엽다,
          오른쪽: 사과는·바나나는·포도는·과일이다·맛있다.
        \item ``사과는''가 과일 군집에 완전히 속지 않고 위쪽으로
          튀어 나온 것 = ``사과는''가 ``맛있다''와 \textbf{두 번} 함께
          나온 반면 바나나는/포도는 한 번 $\to$ \textbf{빈도 차이}
          $\to$ 노름 차이 $\to$ 위치 차이.
      \end{itemize}
      {\footnotesize 말뭉치가 작을수록 이 \textbf{빈도 효과}가 군집
      구조를 어지럽힌다 --- ``문맥이 비슷한 단어가 벡터 공간에서
      가깝다''는 핵심 성질이 이미 작은 스케일에서 확인되는 것.}
    \end{column}
  \end{columns}
\end{frame}
